\documentclass[10pt,a4paper]{amsart}
\usepackage{amsmath,amssymb,amsthm,graphicx}
% Retained preamble for the Stacks Project chapter {CH} (shared project preamble).
% The chapter file pinned-original-source.tex only carries \input{{preamble}} on its first line,
% so no preamble text lives inside the chapter itself.  The Stacks Project uses one single shared
% preamble for every chapter: the snapshot stages it three times (sources/stacks-chow/preamble.tex,
% sources/stacks-etale-cohomology/preamble.tex and focused-seeds/AGP/sources/stacks-snapshot-20261011/
% preamble.tex) and those copies are identical in size.  The block below is copied mechanically,
% line for line, from sources/stacks-chow/preamble.tex lines 14-18 and 154-217: the Xy-pic setup and
% the native theorem declarations together with the mathematical macros.  Nothing was added, removed
% or redefined; the wrapper class and the amsmath/amssymb/amsthm/graphicx base packages are supplied
% by the assembler.
% For commutative diagrams we use Xy-pic
% Chapter: stacks-cohomology.tex -- Cohomology of Sheaves.
\usepackage[all]{xy}

% We use 2cell for 2-commutative diagrams.
\xyoption{2cell}

% Theorem environments.
%
\theoremstyle{plain}
\newtheorem{theorem}[subsection]{Theorem}
\newtheorem{proposition}[subsection]{Proposition}
\newtheorem{lemma}[subsection]{Lemma}

\theoremstyle{definition}
\newtheorem{definition}[subsection]{Definition}
\newtheorem{example}[subsection]{Example}
\newtheorem{exercise}[subsection]{Exercise}
\newtheorem{situation}[subsection]{Situation}

\theoremstyle{remark}
\newtheorem{remark}[subsection]{Remark}
\newtheorem{remarks}[subsection]{Remarks}

\numberwithin{equation}{subsection}

% Macros
%
\def\lim{\mathop{\mathrm{lim}}\nolimits}
\def\colim{\mathop{\mathrm{colim}}\nolimits}
\def\Spec{\mathop{\mathrm{Spec}}}
\def\Hom{\mathop{\mathrm{Hom}}\nolimits}
\def\Ext{\mathop{\mathrm{Ext}}\nolimits}
\def\SheafHom{\mathop{\mathcal{H}\!\mathit{om}}\nolimits}
\def\SheafExt{\mathop{\mathcal{E}\!\mathit{xt}}\nolimits}
\def\Sch{\mathit{Sch}}
\def\Mor{\mathop{\mathrm{Mor}}\nolimits}
\def\Ob{\mathop{\mathrm{Ob}}\nolimits}
\def\Sh{\mathop{\mathit{Sh}}\nolimits}
\def\NL{\mathop{N\!L}\nolimits}
\def\CH{\mathop{\mathrm{CH}}\nolimits}
\def\proetale{{pro\text{-}\acute{e}tale}}
\def\etale{{\acute{e}tale}}
\def\QCoh{\mathit{QCoh}}
\def\Ker{\mathop{\mathrm{Ker}}}
\def\Im{\mathop{\mathrm{Im}}}
\def\Coker{\mathop{\mathrm{Coker}}}
\def\Coim{\mathop{\mathrm{Coim}}}

% Boxtimes
%
\DeclareMathSymbol{\boxtimes}{\mathbin}{AMSa}{"02}

%
% Macros for moduli stacks/spaces
%
\def\QCohstack{\mathcal{QC}\!\mathit{oh}}
\def\Cohstack{\mathcal{C}\!\mathit{oh}}
\def\Spacesstack{\mathcal{S}\!\mathit{paces}}
\def\Quotfunctor{\mathrm{Quot}}
\def\Hilbfunctor{\mathrm{Hilb}}
\def\Curvesstack{\mathcal{C}\!\mathit{urves}}
\def\Polarizedstack{\mathcal{P}\!\mathit{olarized}}
\def\Complexesstack{\mathcal{C}\!\mathit{omplexes}}
% \Pic is the operator that assigns to X its picard group, usage \Pic(X)
% \Picardstack_{X/B} denotes the Picard stack of X over B
% \Picardfunctor_{X/B} denotes the Picard functor of X over B
\def\Pic{\mathop{\mathrm{Pic}}\nolimits}
\def\Picardstack{\mathcal{P}\!\mathit{ic}}
\def\Picardfunctor{\mathrm{Pic}}
\def\Deformationcategory{\mathcal{D}\!\mathit{ef}}
\newenvironment{reference}{}{}
\newenvironment{slogan}{}{}
\newenvironment{history}{}{}
\usepackage{cleveref}
\crefname{originalsection}{original section}{original sections}
\Crefname{originalsection}{Original section}{Original sections}
\makeatletter
\newlabel{sheaves-section-open-immersions}{{009Z}{source}{}{sourcepointer-1}{}}
\newlabel{sheaves-section-open-immersions@cref}{{[originalsection][0][]009Z}{[1][0][]source}{}{}{}}
\makeatother
\crefname{originallemma}{original lemma}{original lemmas}
\Crefname{originallemma}{Original lemma}{Original lemmas}
\makeatletter
\newlabel{sheaves-lemma-j-shriek-modules}{{00A7}{source}{}{sourcepointer-2}{}}
\newlabel{sheaves-lemma-j-shriek-modules@cref}{{[originallemma][0][]00A7}{[1][0][]source}{}{}{}}
\makeatother
\crefname{originallemma}{original lemma}{original lemmas}
\Crefname{originallemma}{Original lemma}{Original lemmas}
\makeatletter
\newlabel{categories-lemma-yoneda}{{001P}{source}{}{sourcepointer-3}{}}
\newlabel{categories-lemma-yoneda@cref}{{[originallemma][0][]001P}{[1][0][]source}{}{}{}}
\makeatother
\crefname{originallemma}{original lemma}{original lemmas}
\Crefname{originallemma}{Original lemma}{Original lemmas}
\makeatletter
\newlabel{homology-lemma-characterize-injectives}{{0136}{source}{}{sourcepointer-4}{}}
\newlabel{homology-lemma-characterize-injectives@cref}{{[originallemma][0][]0136}{[1][0][]source}{}{}{}}
\makeatother
\title{Restriction of an injective module on a ringed space to a smaller open is surjective}
\author{The Stacks Project Authors}
\date{stacks-project:cohomology.tex:local-snapshot-20261011; GNU FDL 1.2 or later}
\begin{document}
\maketitle
\noindent\textit{Editorial note.} All delivered mathematical text is the author's own, copied from the pinned Stacks Project chapter file cohomology.tex; no statement or proof step was rewritten, reordered or removed. The counted statement opens with the author's own slogan 'Injectives are flasque'.\par\medskip
\noindent\textit{Editorial note.} The chapter's retained preamble does not carry the project's no-op definitions of the display environments, so a second verbatim preamble span supplies the project's own no-op definitions of reference, slogan and history, taken from the snapshot at sources/stacks-varieties/retained-preamble.tex lines 75-77 (themselves the project's wrapper definitions). Nothing was invented and no macro was redefined.\par\medskip
\noindent\textit{Editorial note.} All four referenced results live in earlier chapters, so their locators point at the verbatim lines of the project's own tag list staged in the snapshot (sources/stacks-chow/tags.txt) that carry the canonical tags.\par\medskip
\noindent\textit{Allowed standard prerequisites (editorial).}\begin{itemize}
\item Extension by zero of the structure sheaf along an open immersion and its left adjointness to restriction. Sheaves, Section sheaves-section-open-immersions (tag 009Z) and Sheaves, Lemma sheaves-lemma-j-shriek-modules (tag 00A7); invoked at proof lines 728-740. Supply that j\_! O\_U represents the functor F -> F(U) and that restriction is its right adjoint; the Yoneda comparison is retained in the counted span.
\item Yoneda's lemma. Categories, Lemma labelled categories-lemma-yoneda (canonical Stacks tag 001P); invoked at proof line 750. Supplies that the canonical map j'\_! O\_\{U'\} -> j\_! O\_U corresponds to the restriction map on U-sections.
\item Characterisation of injective objects by surjectivity of Hom out of a monomorphism. Homology, Lemma labelled homology-lemma-characterize-injectives (canonical Stacks tag 0136); invoked at proof lines 761-762. Supplies the final step: injectivity of I turns the monomorphism j'\_! O\_\{U'\} -> j\_! O\_U into a surjection on Hom groups, which is the asserted surjectivity of the restriction map.
\end{itemize}\medskip
\ifdefined\hypertarget\hypertarget{sourcepointer-1}{}\fi
\noindent\textit{Source reference.} Original section 009Z: sources/stacks-chow/tags.txt, source lines 371--371. Cross-chapter pointer invoked once at proof lines 729-730 for the extension by zero of O\_U; the adjointness computation and the Yoneda step are retained in the counted span.\par\medskip
\ifdefined\hypertarget\hypertarget{sourcepointer-2}{}\fi
\noindent\textit{Source reference.} Original lemma 00A7: sources/stacks-chow/tags.txt, source lines 379--379. Cross-chapter standard prerequisite invoked once at proof line 740 for the displayed identification Hom(j\_! O\_U, F) = F(U); the comparison of the two representing sheaves is retained in the counted span.\par\medskip
\ifdefined\hypertarget\hypertarget{sourcepointer-3}{}\fi
\noindent\textit{Source reference.} Original lemma 001P: sources/stacks-chow/tags.txt, source lines 70--70. Cross-chapter standard prerequisite invoked once at proof line 750 to identify the canonical map of representing sheaves with the restriction map on sections.\par\medskip
\ifdefined\hypertarget\hypertarget{sourcepointer-4}{}\fi
\noindent\textit{Source reference.} Original lemma 0136: sources/stacks-chow/tags.txt, source lines 1424--1424. Cross-chapter standard prerequisite invoked once at proof lines 761-762 for the criterion turning the monomorphism into a surjection on Hom groups; the final identification with the restriction map is retained in the counted span.\par\medskip
\begin{lemma}
\label{lemma-injective-restriction-surjective}
\begin{slogan}
Injectives are flasque.
\end{slogan}
Let $X$ be a ringed space.
Let $U' \subset U \subset X$ be open subspaces.
For any injective $\mathcal{O}_X$-module $\mathcal{I}$ the
restriction mapping
$\mathcal{I}(U) \to \mathcal{I}(U')$ is surjective.
\end{lemma}

\begin{proof}
Let $j : U \to X$ and $j' : U' \to X$ be the open immersions.
Recall that $j_!\mathcal{O}_U$ is the extension by zero of
$\mathcal{O}_U = \mathcal{O}_X|_U$, see
Sheaves, Section \ref{sheaves-section-open-immersions}.
Since $j_!$ is a left adjoint to restriction we see that
for any sheaf $\mathcal{F}$ of $\mathcal{O}_X$-modules
$$
\Hom_{\mathcal{O}_X}(j_!\mathcal{O}_U, \mathcal{F})
=
\Hom_{\mathcal{O}_U}(\mathcal{O}_U, \mathcal{F}|_U)
=
\mathcal{F}(U)
$$
see Sheaves, Lemma \ref{sheaves-lemma-j-shriek-modules}.
Similarly, the sheaf $j'_!\mathcal{O}_{U'}$ represents the
functor $\mathcal{F} \mapsto \mathcal{F}(U')$.
Moreover there
is an obvious canonical map of $\mathcal{O}_X$-modules
$$
j'_!\mathcal{O}_{U'} \longrightarrow j_!\mathcal{O}_U
$$
which corresponds to the restriction mapping
$\mathcal{F}(U) \to \mathcal{F}(U')$ via Yoneda's lemma
(Categories, Lemma \ref{categories-lemma-yoneda}). By the description
of the stalks of the sheaves
$j'_!\mathcal{O}_{U'}$, $j_!\mathcal{O}_U$
we see that the displayed map above is injective (see lemma cited above).
Hence if $\mathcal{I}$ is an injective $\mathcal{O}_X$-module,
then the map
$$
\Hom_{\mathcal{O}_X}(j_!\mathcal{O}_U, \mathcal{I})
\longrightarrow
\Hom_{\mathcal{O}_X}(j'_!\mathcal{O}_{U'}, \mathcal{I})
$$
is surjective, see
Homology, Lemma \ref{homology-lemma-characterize-injectives}.
Putting everything together we obtain the lemma.
\end{proof}

\end{document}
